% ==================================================================
%  TRANSFER THEORY (pp. 6--7): Pillar 3
% ==================================================================

\section{Transfer Theory}
\label{sec:transfer}

When a practitioner deploys an optimized agent design on a new task
domain, two questions arise.  Can the model selection be reused?  Can the
routing policy be reused?  This section formalizes both questions, proves
they are independent, and derives a three-case decision rule with explicit
thresholds.

\subsection{Two Equivalence Notions}
\label{sec:equivalences}

Let $\tau_S$ (source) and $\tau_T$ (target) be packed-family distributions
over $\Mn$ modes with $m$ models.

\begin{definition}[Post-training equivalence]
\label{def:pt-equiv}
$\tau_S$ and $\tau_T$ are \emph{post-training equivalent} if their
competence profiles coincide:
$p_S(s,M) = p_T(s,M)$ for all modes~$s$ and models~$M$.
\end{definition}

\begin{definition}[Orchestration equivalence]
\label{def:orch-equiv}
$\tau_S$ and $\tau_T$ are \emph{orchestration equivalent} if the optimal
routing policy is the same: $q^*(\tau_S) = q^*(\tau_T)$.
\end{definition}

\begin{proposition}[Independence]
\label{prop:independence}
Post-training equivalence and orchestration equivalence are logically
independent.  Concretely, with $\Mn{=}2$ modes and $m{=}2$ models:
(a)~There exist $\tau_A, \tau_B$ that are post-training equivalent but
  not orchestration equivalent.
(b)~There exist $\tau_C, \tau_D$ that are orchestration equivalent but
  not post-training equivalent.
\end{proposition}

\begin{proof}[Proof sketch]
For~(a), take identical competence matrices but
$\pi_A = (0.7, 0.3)$, $\pi_B = (0.3, 0.7)$.
The posterior routing distributions differ
($\KL(\pi_1^A \| \pi_1^B) \approx 0.32$) because mode frequencies
shift, even though per-mode success probabilities are identical.
For~(b), take $\pi_C = \pi_D = (0.5, 0.5)$ with the same model-to-mode
assignment ordering but different magnitudes
($p_C(1,M_{\mathrm{big}}) = 0.9$ vs.\
$p_D(1,M_{\mathrm{big}}) = 0.6$).
The routing policy is identical (same ordering), but the competence
profiles differ.
Full constructions in Appendix~\ref{app:independence}.
\end{proof}

The independence is not a technicality: it means the transfer decision
space is genuinely two-dimensional.  Knowing that the models perform
similarly (post-training) tells you nothing about whether the routing
structure transfers.

\subsection{Transfer Divergences and the Three-Case Rule}
\label{sec:transfer-rule}

We measure source--target distance with two divergences:
\begin{align}
  \alpha_{\mathrm{PT}} &= \max_{s,M}\,
    \bigl|\log p_S(s,M) - \log p_T(s,M)\bigr|
    \quad \text{(competence shift)},
    \label{eq:alpha-pt} \\
  \alpha_O &= \KL(\pi_S \| \pi_T)
    \quad \text{(orchestration shift)}.
    \label{eq:alpha-o}
\end{align}

\begin{theorem}[Three-case transfer bounds]
\label{thm:transfer}
Let $d_S^*$ be the source-optimal design and $d_T^*$ the target-optimal.
The transfer error
$\Delta = f(d_S^*;\tau_T) - f(d_T^*;\tau_T)$ satisfies:

\medskip
\noindent\textbf{Case 1 (Plug-and-play):}
$\alpha_{\mathrm{PT}} \le \tau_1$, $\alpha_O \le \tau_2$.
\begin{equation}
  \Delta \;\le\; 2\,\alpha_{\mathrm{PT}} + \sqrt{2\,\alpha_O}\cdot\log\Mn.
\end{equation}

\noindent\textbf{Case 2 (Re-route):}
$\alpha_{\mathrm{PT}} \le \tau_1$, $\alpha_O > \tau_2$.
Keep $M_S^*$, re-estimate routing from $N_{\mathrm{re}}$ target samples:
\begin{equation}
  \Delta \;\le\; \alpha_{\mathrm{PT}} + 2\dacc, \quad
  N_{\mathrm{re}} = O\!\Bigl(\frac{\Mn\log(\Mn/\dconf)}
    {\dacc^2\,\pmin^2}\Bigr).
\end{equation}

\noindent\textbf{Case 3 (Start over):}
$\alpha_{\mathrm{PT}} > \tau_1$.
Run full Algorithm~1 on~$\tau_T$.
$\Delta \le 2\dacc$ with $N_{\mathrm{full}} = O(m\cdot\Mn/R^2)$.
\end{theorem}

\begin{proof}[Proof sketch]
Decompose the transfer error term-by-term using the four-term identity.
The cost term is distribution-independent.  The competence shift is
bounded by $2\alpha_{\mathrm{PT}}$ via the definition of the max-log
divergence.  The information-gain and routing-mismatch shifts are bounded
by $\sqrt{2\alpha_O}\cdot\log\Mn$ via Pinsker's inequality and
entropy continuity.  Cases~2 and~3 follow by composing the transfer bound
with the near-optimality guarantee (Proposition~\ref{prop:near-opt}).
Full proof in Appendix~\ref{app:transfer}.
\end{proof}

The thresholds are derived by inverting the Case~1 bound to a target
accuracy~$\dtarget$:

\bigskip
\noindent\fbox{\parbox{0.95\textwidth}{%
\textbf{Decision Rule: Three-Case Transfer Protocol}

\medskip
\textbf{Input:} Source design $d_S^*$, divergences $\alpha_{\mathrm{PT}},
\alpha_O$, target accuracy $\dtarget$.

\textbf{Thresholds:}
$\tau_1 = \dtarget/2$, \quad
$\tau_2 = \dtarget^2/(2\log^2\Mn)$.

\textbf{Decision:}
\begin{enumerate}
\item $\alpha_{\mathrm{PT}} \le \tau_1$ \textbf{and}
  $\alpha_O \le \tau_2$:
  \textbf{Plug-and-play.} Transfer directly.
  Cost: 0 new evaluations.
\item $\alpha_{\mathrm{PT}} \le \tau_1$ \textbf{but}
  $\alpha_O > \tau_2$:
  \textbf{Re-route.} Keep model, re-estimate
  routing.  Cost: $O(\Mn/(\dacc^2\pmin^2))$ evaluations
  (factor $m$ cheaper than full).
\item $\alpha_{\mathrm{PT}} > \tau_1$:
  \textbf{Start over.} Run full
  Algorithm~1 on target.
\end{enumerate}
}}
\bigskip

\subsection{Transfer Savings}

\begin{proposition}[Transfer savings]
\label{prop:transfer-savings}
In Case~1, the target cost is zero (amortized over the source).
In Case~2, the pilot cost drops by a factor of~$m$ relative to
Case~3, because only one model (the transferred $M_S^*$) requires
re-piloting.  In Case~3, no savings accrue from transfer.
\end{proposition}

\subsection{Transfer Experiment: EXP12}
\label{sec:exp12}

We sweep a $7 \times 6$ grid of
$(\alpha_{\mathrm{PT}}, \alpha_O)$ values,
generating target distributions at controlled divergence from a source
with $\Mn{=}5$ modes and $m{=}5$ models.  At each grid point, we measure
transfer regret (100 Monte Carlo replications) and record the decision
rule's case assignment.

The decision rule correctly assigns the transfer case in 73.8\% of
borderline cells (those within $2\times$ of a threshold).
On interior cells (far from thresholds), accuracy exceeds 95\%.
The errors are one-directional: the rule is conservative, defaulting to
more expensive strategies when uncertain.  This conservatism is by
design---the thresholds use worst-case bounds.

Monte Carlo verification of both independence directions (EXP11, 200 replications, zero violations) is in Appendix~\ref{app:exp11}.
